% Section 6: Conclusion (400 words target)
% TODO: Complete after results available

We have conducted the first systematic empirical test of the Modifiable Areal Unit Problem in algorithmic redistricting by comparing recursive bisection across three spatial resolutions: census tracts, block groups, and census blocks. [Summary of key findings TBD].

\paragraph{Contribution.} This work establishes the generalizability---or lack thereof---of automated redistricting findings to different spatial scales. [If robust: Our results validate census tracts as appropriate atomic units for redistricting research, enabling confident extrapolation of computational findings to policy contexts.] [If sensitive: Our results demonstrate that spatial resolution is a critical methodological choice with substantive implications for district design and VRA compliance analysis.]

\paragraph{Implications.} For researchers, this work [establishes/challenges] the robustness of algorithmic redistricting methods. For practitioners, it provides [validation/guidance] on resolution choice in automated district design. For courts, it [confirms/complicates] the use of computational evidence in redistricting litigation.

\paragraph{Final statement.} The choice of spatial resolution is not merely a technical detail---it is a methodological decision with consequences for democratic representation.
